  \section{System}
  \label{sec:system}

  \subsection{The sources}
  \label{sec:sources}

  Building the advisory starts at the bottom of \Cref{fig:architecture}, with the three sources
  every later number is drawn from (\Cref{tab:sources}). They fix what the layers above can say,
  and each can be trusted to a different degree.

  \begin{table}[htbp]
      \centering
      \begin{tabular}{@{}lll@{}}
      \toprule
      \textbf{Source} & \textbf{Provides} & \textbf{Standing} \\
      \midrule
      DT models & thermal and control behaviour & calibrated against the box \\
      Domain knowledge & tempeh and heater limits & sourced or estimated \\
      Current state & the running batch's condition & as published by the twin \\
      \bottomrule
      \end{tabular}
      \caption{Three sources, what each contributes, and where its authority comes from.}
      \label{tab:sources}
  \end{table}

  \textbf{The DT's models.} Four modules are taken from the twin: two plants, one
  carrying the air temperature and the box's heat loss and one adding the heater element, the
  controller's switching logic, and the wrapper that makes that logic a model the solver can run
  \cite{gomes2025incubator}. They are used exactly as the twin defines them. Their calibrated
  parameters arrive with the current state rather than being written into this code, so a
  recalibration on the twin's side reaches the advisory with nothing here edited.

  The parameters were fitted
  with the box empty and the lid closed, and room temperature is held constant over the horizon.

  \textbf{The current state.} A single document stands in for what the twin publishes about the
  running batch: sensor readings, actuator settings, the filtered estimate of the box's
  temperature, the controller's settings, the calibrated model parameters, the open-lid parameters,
  and how far the batch has progressed. It is the only source that changes between one question
  and the next.

  \textbf{Domain knowledge.} One document holds what the twin does not (\Cref{tab:domain}). The
  temperature values follow the tempeh literature
  \cite{steinkraus1960tempeh,tempehinfo_temperature}, the bands, budgets and heater limits are
  reasoned values adopted for this system and marked as estimated. None has been calibrated
  against a batch that actually spoiled.

  \begin{table}[htbp]
    \centering
    \begin{tabular}{@{}llc@{}}
      \toprule
      \textbf{Entry} & \textbf{Value} & \textbf{Basis} \\
      \midrule
      \multicolumn{3}{@{}l}{\textit{Fermentation phases, over a 24 hour batch}} \\
      \quad Lag            & 0 to 11 h, cold tolerance $\times 0.5$  & sourced \\
      \quad Active growth  & 11 to 21 h, $\times 1.75$               & sourced \\
      \quad Maturation     & 21 to 24 h, $\times 1.15$               & sourced \\
      \midrule
      \multicolumn{3}{@{}l}{\textit{Culture temperature bounds}} \\
      \quad Optimal range  & 30 to 37 \si{\celsius}, best near 34.5  & sourced \\
      \quad Growth limits  & 22 to 41 \si{\celsius}                  & sourced \\
      \quad Damage ceiling & 42 \si{\celsius}                        & sourced \\
      \midrule
      \multicolumn{3}{@{}l}{\textit{Safe air band, by phase}} \\
      \quad Lag            & 27 to 35 \si{\celsius}                  & estimated \\
      \quad Active growth  & 23 to 39 \si{\celsius}                  & estimated \\
      \quad Maturation     & 25 to 37 \si{\celsius}                  & estimated \\
      \midrule
      \multicolumn{3}{@{}l}{\textit{Time a batch may spend below its band}} \\
      \quad Still safe     & 600 s                                   & estimated \\
      \quad Recoverable    & 900 s                                   & estimated \\
      \midrule
      \multicolumn{3}{@{}l}{\textit{Heater limits}} \\
      \quad Element ceiling      & 60 \si{\celsius}                  & estimated \\
      \quad Longest run, then rest & 1800 s, then 300 s              & estimated \\
      \quad Share of any hour on & 0.75                              & estimated \\
      \bottomrule
    \end{tabular}
    \caption{Domain knowledge overview.}
    \label{tab:domain}
  \end{table}



  \subsection{Tool helpers}
  \label{sec:helpers}

  Two small modules stand between the sources and the simulation. Neither decides anything. They
  make the sources available in a form the simulation can use.

  \textbf{Reading the state.} The first opens the state document and validates it. It is the only
  place that document is read, so a state that is malformed, incomplete or internally inconsistent
  fails here with a clear message rather than reaching a simulation and producing a plausible
  wrong answer.

  \textbf{Reading the domain knowledge.} The second loads the thresholds in \Cref{tab:domain} and
  answers two questions: which phase of fermentation the batch has reached, derived from how long
  it has been running since an incubator has no sensor for it, and which temperature band counts
  as safe in that phase. Both are needed before a trajectory means anything, which is why they are
  settled here rather than downstream.

  \subsection{The simulation tool}
  \label{sec:simulation}

  The simulation is the capability the language layer uses to answer a what-if question about the
  lid: given an opening, it predicts what the box will do, starting from the state it is in now.

  \textbf{Inputs.} The language layer supplies only the opening: how long the lid stays open and
  when it opens, by default immediately. Everything else is read from the current state, including
  the starting temperatures, the room temperature, the controller settings, the calibrated
  parameters and the open-lid disturbance. The skills have the language layer choose which opening
  to simulate, not the conditions it starts from, although the tool can be pointed at another
  state file, a route \Cref{sec:runs} returns to. An opening of zero length gives the undisturbed baseline,
  so the same tool also answers what happens if the lid stays shut. A negative duration, an opening
  longer than six hours or a state that fails validation is refused with a plain message, and
  nothing is replaced by a default.

  \textbf{What it does.} It runs the twin's own plant and controller together in closed loop
  \cite{gomes2025incubator}, so the heater reacts to the opening as it would on the box. An open
  lid is represented by multiplying the box's heat-loss coefficient, for as long as the lid is
  open, by a ratio taken from the twin's more detailed calibration. The twin declares that
  coefficient as an input rather than a fixed parameter, so supplying a value that changes over
  time is ordinary use of its interface, and none of its equations, parameters or code is changed.
  The run continues for an hour after the lid closes, long enough to observe recovery.

  \textbf{What it returns.} A result has three parts. The summary holds the numbers an answer is
  built from: the lowest air temperature and when it occurs, the total drop, the peak heater
  temperature, the time to recover, and the time spent outside the safe band. The trajectory is
  thinned to about a hundred points, which is the reduction step described in
  \Cref{sec:related-dt}, every number is still computed on the full solver grid. The inputs record
  lists what was actually used, including the fermentation phase and the safe band that applied.
  The tool computes and does not interpret: whether any of this is acceptable for the batch is
  decided by the risk assessment (\Cref{sec:analysis}).

  \textbf{Why this suits the language layer.} Three properties make the tool usable by a model
  that must not compute on its own:

  \begin{itemize}
    \item \textbf{It is phrased in the operator's terms.} A duration and a start time are all a call
      needs, so turning a question into a call requires no knowledge of thermal parameters. This
      removes the first obstacle in \Cref{sec:gap}.
    \item \textbf{Every number carries its assumptions.} Because the inputs record is part of the
      result, an explanation can state what a prediction rests on, such as the phase the batch was
      judged to be in, without the model reconstructing it.
    \item \textbf{It follows the twin.} Parameters and controller settings arrive with the state, so
      after the twin recalibrates or re-tunes its controller (\Cref{sec:incubator-dt}), the next
      prediction uses the new values with nothing here changed. And because an opening is only a
      duration and a start time, a question can be explored by comparison, such as opening now
      against opening later, the same pattern of repeated calls the heater-boost advisory builds
      on (\Cref{sec:architecture}).
  \end{itemize}

  % These three properties are what the discussion (sec:flexibility) tests.

  \subsection{The risk assessment}
  \label{sec:analysis}

  The risk assessment turns the simulation's numbers into a verdict about the batch, using rules in
  ordinary code so that the thresholds behind a verdict are stated rather than inferred. From the
  batch's age when the lid opens it takes the phase, so a later opening is judged in the phase it
  falls in, and from \Cref{tab:domain} that phase's safe air band,
  its cold tolerance, and the two exposure budgets that the tolerance scales, so a young batch is
  treated as less forgiving than one generating its own heat.

  Time below the band within the first budget is safe, between the two budgets recoverable, and
  beyond the second high risk, time above the band but under the damage ceiling of
  42 \si{\celsius} turns an otherwise safe verdict into recoverable. Air above the ceiling, which
  an open lid cannot cause but the heater boost could, and a run that never returns to the
  setpoint are both high risk, the second because recovery cannot be confirmed, not because harm
  was shown.

  A recoverable verdict, the only one where what the operator does next changes the outcome, comes
  with the time the lid must then stay shut. Every verdict carries the number that decided it and
  the limits it inherits from the models and thresholds behind it (\Cref{sec:sources}).

  For example, a five-minute opening in active growth, with the air at 30.98 \si{\celsius} and
  the room at 22.4 \si{\celsius}, is predicted to take the air down to 23.81 \si{\celsius}. That
  stays inside the band, so the verdict is safe, but with a margin of only 0.8 \si{\celsius}.


  \subsection{The lid-opening advisory}
  \label{sec:skill}

  The lid-opening advisory, like each advisory in the language layer of \Cref{fig:architecture},
  is written as instructions rather than as code: a short document, a skill
  (\Cref{sec:architecture}), loaded when a question is about opening the lid. It names the single
  tool the advisory may call and fixes the order in which the work is done, so no orchestration
  program of our own sits between the operator and the tool.

  What the document fixes is how far the model's discretion reaches. The model may recognise that
  a question concerns the lid, read the opening as a duration, and choose the words, it may not
  fill in what the operator left out, and a question the tool will not accept returns to the
  operator instead of being completed by the model. The rest is relayed: every number in an answer
  must appear in the tool's output. The answer is short and in a fixed order: the verdict with the
  reason that decided it, what happens to the air, and for a recoverable opening how long to keep
  the lid shut. It ends by offering the limitations, the caveats and the standing of the thresholds
  (\Cref{sec:sources}) are explained only when asked for, and then never softened.

  \subsection{The heater-boost advisory}
  \label{sec:boost}

  When an opening is judged high risk, the operator can ask whether warming the box first would make
  it acceptable. The only setting that can be changed on a running incubator is the controller's
  setpoint, so a boost is a plan: raise the setpoint now, open the lid once the air has reached it,
  and set it back when the lid closes. This is the twin's own example of heating ahead of a power
  supply replacement (\Cref{sec:incubator-dt}), applied to the lid.

  The tool tries whole degrees above the current setpoint, up to the top of the phase's safe band.
  Each candidate is simulated as one continuous run, through the pre-heat, the opening and the
  recovery, and judged by the unchanged risk assessment (\Cref{sec:analysis}), so a boosted opening
  is measured and judged exactly as a plain one. A candidate is accepted if it is judged
  recoverable or safe, stays within the heater limits of \Cref{tab:domain}, and clears the
  exposure budget by more than the interval between the solver's samples, since a smaller margin
  can flip with a fraction of a second more or less pre-heat. The tool returns the lowest accepted
  setpoint as a plan, or the reason none exists.

  In this model a boost achieves little. The air loses its extra heat within seconds of the lid
  opening, so the gain comes only from heat stored in the heating element, which shortens recovery
  once the lid is shut. In the lag phase with the room at 17.6 \si{\celsius}, for example, a
  five-minute opening leaves the air below its band for 488.7 s against a budget of 450 s and is
  high risk. A boost to 34 \si{\celsius}, with the lid opened after 351.0 s, brings this to
  441.0 s and the verdict to recoverable. A boost to 33 \si{\celsius} is also judged recoverable,
  but clears the budget by 3.0 s against a sample interval of 3.1 s and is rejected.

  The advisory is a second skill, separate from the first, so it was added on top of the same
  tools without changing the twin or the lid-opening advisory, how far this carries is discussed
  in \Cref{sec:generality}. Its rules follow \Cref{sec:skill},
  with two additions: the plan is passed on exactly as returned, together with how narrowly it
  clears the budget, and when no boost works a shorter opening is suggested but not searched for,
  since a search run by the model would be a numerical decision taken in the language layer.
